% TOP_PROBLEM: {"id": 3209, "title": "Weak pentagon problem", "category_id": 3, "category": {"id": 3, "name": "graph_theory", "display_name": "Graph Theory", "description": "Problems involving graphs, networks, and their properties.", "slug": "graph-theory", "order_index": 3, "created_at": "2026-07-31T15:26:25.671Z"}, "status": "open"}
% FIRST_POSTED: null
% ENTRY_KIND: statement_only
% Imported from: batch02.tex; UnsolvedMath ID 3209
% Compile twice: lualatex 3209.tex
\documentclass[11pt]{article}
\usepackage{fontspec}
\setmainfont{Latin Modern Roman}
\usepackage{amsmath,amssymb,mathtools,unicode-math}
\setmathfont{Latin Modern Math}
\usepackage{xurl,hyperref}
\hypersetup{hidelinks,pdftitle={UnsolvedMath Problem 3209}}
\newcommand{\sref}[2]{\hyperlink{source:#1:#2}{[S#2]}}
\newcommand{\cataloguescope}[1]{\paragraph{Mathematical question.}#1}
\begin{document}
% UnsolvedMath ID 3209; source catalogue code OPG-434
\setcounter{section}{3208}
\section{Weak pentagon problem}\label{3209}
{\small\sffamily UnsolvedMath ID 3209 · Graph Theory}
\subsection{Definitions and mathematical statement}

\paragraph{Shared notation.}$\mathbb N=\{1,2,\ldots\}$, and $\mathbb Z,\mathbb Q,\mathbb R,\mathbb C$ denote the integers, rationals, reals and complex numbers. Any other conventions are specified locally.


A cubic graph is a finite simple $3$-regular graph; its girth is the length of a shortest cycle. For $j\ge2$, the projective cube $PQ_j$ is the quotient of the $(j+1)$-dimensional cube (binary strings, adjacent when differing in one coordinate) by identification of antipodal strings. Thus $PQ_4$ is the Clebsch graph. A homomorphism sends adjacent vertices to adjacent vertices. A cut-continuous edge map $f:E(G)\to E(H)$ has the property that $f^{-1}(\delta_H(X))$ is a cut of $G$ for every vertex subset $X$ of $H$.
\paragraph{Statement recorded in the supplied source.}
Does every triangle-free cubic graph admit an edge coloring with five colors such that deleting any one color class leaves a bipartite graph? Equivalently, does it have a homomorphism to $PQ_4$, or a cut-continuous map to $C_5$? Also determine the largest integer $k\ge1$ for which there is $g_k$ such that every cubic graph of girth at least $g_k$ maps homomorphically to $PQ_{2k}$. The latter condition is equivalent to an edge coloring with $2k+1$ colors whose individual color-class complements are bipartite, or to a cut-continuous map to $C_{2k+1}$. These edge colorings are not required to be proper. \sref{3209}{2}
\subsection{Short English statement}
Can five color classes each meet every odd cycle in a triangle-free cubic graph, and how far can the analogous high-girth projective-cube statement be extended?
\subsection{Sources}
\begin{description}
\item[\hypertarget{source:3209:1}{S1}] ULAM.AI and the UnsolvedMath Contributors, UnsolvedMath: A Curated Collection of Open Mathematics Problems, record 3209 (OPG-434). CC BY 4.0. \url{https://huggingface.co/datasets/ulamai/UnsolvedMath}.
\item[\hypertarget{source:3209:2}{S2}] Open Problem Garden, Weak pentagon problem, original node 434, statement and mathematical discussion. \url{https://www.openproblemgarden.org/op/weak_pentagon_problem}.
\end{description}
\label{end:3209}
\end{document}
